\documentclass[11pt, a4paper]{amsart}
\usepackage{amsmath, amssymb, amsthm, amsfonts}
\usepackage{mathtools}
\usepackage[margin=1in]{geometry}
\usepackage{hyperref}
\usepackage{tikz-cd}

\newtheorem{theorem}{Theorem}[section]
\newtheorem{lemma}[theorem]{Lemma}
\newtheorem{proposition}[theorem]{Proposition}
\newtheorem{corollary}[theorem]{Corollary}
\newtheorem{fact}[theorem]{Standard Fact}
\theoremstyle{definition}
\newtheorem{definition}[theorem]{Definition}
\newtheorem{remark}[theorem]{Remark}
\newtheorem{notation}[theorem]{Notation}
\newtheorem{example}[theorem]{Example}

\newcommand{\R}{\mathbb{R}}
\newcommand{\Q}{\mathbb{Q}}
\newcommand{\C}{\mathbb{C}}
\newcommand{\Z}{\mathbb{Z}}
\newcommand{\N}{\mathbb{N}}
\newcommand{\HH}{\mathbb{H}}
\newcommand{\OO}{\mathbb{O}}
\newcommand{\id}{\operatorname{id}}
\newcommand{\im}{\operatorname{im}}
\newcommand{\tr}{\operatorname{tr}}
\newcommand{\spec}{\operatorname{spec}}
\newcommand{\op}{\operatorname{op}}
\newcommand{\Cl}{\mathrm{Cl}}
\newcommand{\inv}{\mathrm{inv}}
\newcommand{\cyl}{\mathrm{cyl}}
\renewcommand{\Im}{\mathrm{Im}}

\title{Cascade--Classical Correspondence IV:\\
  Discrete Closure Dynamics}
\author{}
\date{}

\begin{document}

\maketitle

\begin{abstract}
This paper puts dynamics on the refinement-spaces substrate of paper
P3. At each refinement level $n \geq 0$ and branch $\bullet \in
\{H_4, D_4\}$, we define an explicit quadratic closure functional
\[
  F_n = \alpha R_n + \beta E_n - \gamma Q_n
    = \tfrac{1}{2} \langle \cdot, L_n \cdot \rangle
\]
on the finite-level Hilbert space $X_n := X_n^{0, \bullet} \oplus
X_n^{1, \bullet}$, where $R_n, E_n, Q_n$ are the curvature, kinetic,
and mass quadratic forms built from the discrete coboundary,
graph Laplacian, and mass operator of \S\ref{sec:operators}. The
gradient of $F_n$ is represented by a bounded self-adjoint
operator $L_n \in B(X_n)$
(Theorem~\ref{thm:gradient-operator}); the continuous-time gradient
flow $\dot\Phi_n = -L_n \Phi_n$ has unique global solutions
$\Phi_n(t) = e^{-t L_n} \Phi_n(0)$
(Theorem~\ref{thm:flow-exists}); energy dissipates along the flow
as $d_t F_n(\Phi_n(t)) = -\|L_n \Phi_n(t)\|^2 \leq 0$
(Proposition~\ref{prop:energy-dissipation}).

A discrete-time Euler scheme $\Phi_{n, k+1} = (I - \eta L_n)
\Phi_{n, k}$ is analysed in
Section~\ref{sec:discrete-time}. Energy monotonicity under
small step size is separated cleanly from norm contractivity
(Theorem~\ref{thm:coercive-contraction}), which holds only on
coercive invariant subspaces $Y_n$ of $X_n$ where $L_n$ has a
positive spectral gap.

Refinement compatibility of the full gradient operator
$L_n = \alpha A_n + \beta B_n - \gamma C_n$ is established here
only in the mass-only case $\alpha = \beta = 0$
(Theorems~\ref{thm:refinement-compat} and
\ref{thm:flow-intertwining}); the two constituent results that
\emph{are} proved are the factor-$\tfrac{1}{2}$ coboundary
relation $p^1 d_{n+1} = \tfrac{1}{2} d_n p^0$
(Proposition~\ref{prop:coboundary-refinement}) and the exact
mass-block compatibility $p_{n+1, n} C_{n+1} = C_n p_{n+1, n}$
(Proposition~\ref{prop:adjoint-refinement}). For arbitrary $(\alpha,
\beta, \gamma)$ this paper makes no claim about refinement-
intertwining of $L_n$ (neither positive nor negative); see
Remark~\ref{rmk:refinement-compat-scope}. Downstream papers
needing refinement-intertwining of cascade dynamics in that
regime must supply their own analysis.

The false global theorem ``all correspondences are Borel /
equivariant'' from an earlier foundations draft is replaced by an
explicit generator-family list of \emph{admissible finite-level
intertwiners} (Definition~\ref{def:admissible}), with printed
operator-norm bounds and continuity/Borel measurability following
directly from finite-dimensional functional calculus
(Theorem~\ref{thm:intertwiner-bounds}). This is the fourth of
thirteen planned cascade infrastructure papers.

\emph{No continuum-limit, Gromov--Hausdorff, or
manifold-identification claim is made here. No target-side
metric, hydrodynamic, gauge, spectral, Hecke, cohomology, or CTM
correspondence theorem is in scope; those are P5--P13.}
\end{abstract}

\tableofcontents

%=====================================================================
\section{Introduction and Source Audit}
\label{sec:intro}
%=====================================================================

\subsection{The repair target}

The cascade framework has long referred to ``closure dynamics'' in
heuristic language: a closure functional $\mathcal{F} = \alpha R +
\beta E - \gamma Q$ whose critical points or gradient flows
organise the cascade rungs. The earlier draft of the cascade
correspondence foundations manuscript invoked this at an abstract
global level via a ``Theorem'' asserting that every correspondence
arising from closure dynamics is Borel and symmetry-equivariant
(iter-1 review of that manuscript, Theorem 4.3 of its §4). The
review identified that claim as false in the form stated, and
recommended replacement by case-by-case operator-level bounds.

This paper executes that replacement at the finite-level stage.
We work exclusively on the finite-dimensional Hilbert spaces
$X_n := X_n^{0, \bullet} \oplus X_n^{1, \bullet}$ of paper P3,
define the closure functional as an explicit quadratic form,
identify its gradient as a bounded self-adjoint operator $L_n$,
and derive its dynamics rigorously. Every operator appearing in
downstream constructions (bonding maps, flow maps, symmetry
actions, sector projections, coboundaries, adjoints, Laplacians,
mass operators, and the specific compositions used by P5--P13)
is given an explicit norm bound.

\subsection{What this paper proves}

\begin{enumerate}
\item The closure functional $F_n = \alpha R_n + \beta E_n -
  \gamma Q_n$ is a continuous symmetric quadratic form on $X_n$,
  equivalently $F_n(\Phi) = \tfrac{1}{2}\langle \Phi, L_n \Phi
  \rangle$ for a unique bounded self-adjoint operator $L_n$
  (\S\ref{sec:quadratic}, Theorem~\ref{thm:gradient-operator}).
\item The continuous-time gradient flow $\dot\Phi_n = -L_n \Phi_n$
  has unique global solutions on $[0, \infty)$ given by
  $\Phi_n(t) = e^{-t L_n} \Phi_n(0)$, a finite-dimensional
  $C^0$-semigroup on $X_n$
  (\S\ref{sec:continuous-flow}, Theorem~\ref{thm:flow-exists}).
\item Energy dissipation along the flow:
  $d_t F_n(\Phi_n(t)) = -\|L_n \Phi_n(t)\|^2 \leq 0$
  (\S\ref{sec:energy}, Proposition~\ref{prop:energy-dissipation}).
\item Discrete-time Euler iteration: energy-decreasing under a
  step-size bound; strict norm-contraction only on coercive
  invariant subspaces of $L_n$
  (\S\ref{sec:discrete-time}, Propositions~\ref{prop:euler-energy-decrease}
  and Theorem~\ref{thm:coercive-contraction}).
\item Refinement compatibility for the mass-only case
  ($\alpha = \beta = 0$): the identity $p_{n+1, n} L_{n+1} =
  L_n p_{n+1, n}$ and the flow intertwining
  $p_{n+1, n} e^{-t L_{n+1}} = e^{-t L_n} p_{n+1, n}$ hold
  (\S\ref{sec:refinement-compat},
  Theorems~\ref{thm:refinement-compat} and
  \ref{thm:flow-intertwining}). The constituent intermediate
  identities proved are the factor-$\tfrac{1}{2}$ coboundary
  relation (Proposition~\ref{prop:coboundary-refinement}) and the
  exact mass-block compatibility
  (Proposition~\ref{prop:adjoint-refinement}). For arbitrary
  $(\alpha, \beta, \gamma)$ this paper makes no claim about
  refinement-intertwining of $L_n$; see
  Remark~\ref{rmk:refinement-compat-scope}.
\item Admissible finite-level intertwiners (a finite list of
  generators: bonding maps, cylindrical embeddings, sector
  projections, symmetry actions, coboundaries, adjoints,
  Laplacians, $L_n$, $e^{-tL_n}$, and fibre maps from P2) have
  explicit operator-norm bounds; all admissible intertwiners are
  bounded, continuous, and Borel
  (\S\ref{sec:intertwiners}, Theorem~\ref{thm:intertwiner-bounds}).
\end{enumerate}

\subsection{What this paper does not prove}

To forestall any reader from treating P4 as a continuum-dynamics
or branch-specific-correspondence paper, we record the exclusions
explicitly:

\begin{enumerate}
\item No continuum-limit statement: no Gromov--Hausdorff or
  spectral-convergence claim about the $n \to \infty$ behaviour
  of $F_n, L_n, e^{-tL_n}$, or the bonding maps. Those belong to
  paper P6.
\item No continuum PDE identification. This paper gives no
  heat-equation, Ricci-flow, Navier--Stokes, or Yang--Mills-type
  interpretation of $\dot\Phi_n = -L_n \Phi_n$. Those are P5, P7,
  and P9.
\item No target-side metric, hydrodynamic, gauge, spectral,
  Hecke, cohomology, or CTM/TM correspondence theorem.
\item No universal ``all correspondences are Borel and
  equivariant'' global theorem. Instead
  Section~\ref{sec:intertwiners} gives an explicit generator-family list
  with explicit bounds for the operators later papers actually
  invoke.
\item No claim that $F_n$ has a unique minimiser or that the flow
  converges to the ground state without further spectral
  hypotheses. Energy monotonicity and norm contractivity are
  separated theorems with separate hypotheses.
\item No claim about uniqueness or classification of the closure
  functional at the level of abstract $W^\bullet$-symmetric
  quadratic forms. (The cascade-foundations note F2 contains such
  a claim; P4 does not rely on or reproduce it.)
\end{enumerate}

\subsection{Source audit}

Theorem-grade inputs to this paper:
\begin{description}
\item[\textbf{Internal cascade infrastructure.}]
  Paper P2 (cascade algebraic substrate
  \cite{AlgebraicSubstrate}) for the sector fibres $V_{H_4},
  V_{D_4}, V_{16}, V_8, \Im(\OO)$ and their fixed trace forms.
  Paper P3 (cascade refinement spaces
  \cite{RefinementSpaces}) for the finite-level Hilbert spaces
  $X_n^{0, \bullet}, X_n^{1, \bullet}$, the downward bonding
  maps $p_{n+1, n}$, the adjoint cylindrical embeddings
  $j_{n, n+1}$, and the finite-level Coxeter actions
  (P3 only constructs the vertex inverse-limit Hilbert space and
  defines $\sigma$ on the mixed $\Q(\varphi)/\Q$ vertex form,
  not on the real Hilbert space; we use $\sigma$ here only at
  the rational-form level).
\item[\textbf{Classical mathematics.}]
  Spectral graph theory and finite-graph Hodge theory
  \cite{Chung}; bounded self-adjoint operators and spectral
  calculus in finite dimension \cite{ReedSimonI, HornJohnson};
  linear ODEs with bounded coefficients
  \cite{CoddingtonLevinson, Hartman}.
\end{description}

We do not use as theorem-grade input any cascade-foundations F2
uniqueness claim, any iter-1 foundations.tex Theorem 4.3 global
theorem, or any continuum-limit cascade note.

\subsection{Organization}

Section~\ref{sec:operators} fixes the finite-level discrete
operators (coboundary, Laplacian, mass) and proves their basic
bounds. Section~\ref{sec:quadratic} defines the closure functional
and identifies its gradient. Section~\ref{sec:continuous-flow}
constructs the continuous-time flow. Section~\ref{sec:energy}
derives the energy dissipation identity.
Section~\ref{sec:discrete-time} analyses the discrete-time Euler
scheme and separates energy monotonicity from norm contractivity.
Section~\ref{sec:refinement-compat} establishes refinement
intertwining for the coboundary block (factor $\tfrac{1}{2}$)
and a mass-only refinement intertwining for $L_n$ in the
specialization $\alpha = \beta = 0$, and records the obstruction
to claiming more in the general $\alpha, \beta, \gamma > 0$
regime. Section~\ref{sec:intertwiners} replaces the false
global Borel/equivariance theorem by explicit generator-family
bounds. Section~\ref{sec:handoff} states the citation contract for
downstream papers. Appendix~\ref{app:notation} is a notation
concordance.

%=====================================================================
\section{Finite-Level State Spaces and Discrete Operators}
\label{sec:operators}
%=====================================================================

\subsection{State spaces}

\begin{definition}[Combined state space]
\label{def:Xn}
For $n \geq 0$ and $\bullet \in \{H_4, D_4\}$, let
\[
  X_n := X_n^{0, \bullet} \oplus X_n^{1, \bullet}
\]
denote the direct sum of the vertex-sector and edge-sector
$\ell^2$-spaces from
\cite[Definitions~\texttt{def:Xn-0} and~\texttt{def:Xn-1}]{RefinementSpaces}, with
the direct-sum inner product
$\langle (f, A), (g, B) \rangle_{X_n} := \langle f, g
\rangle_{X_n^0} + \langle A, B \rangle_{X_n^1}$.
The \emph{domain} is $D_n := X_n$; since $X_n$ is
finite-dimensional (Proposition~\ref{prop:Xn-finite}), the domain
is automatically dense in itself and every operator introduced
below is bounded.
\end{definition}

\begin{proposition}[$X_n$ is finite-dimensional]
\label{prop:Xn-finite}
$X_n$ is a finite-dimensional real Hilbert space with
$\dim_\R X_n = \dim_\R X_n^{0, \bullet} + \dim_\R X_n^{1, \bullet}
< \infty$. The dimension bounds are given by
\cite[Proposition~\texttt{prop:Xn-Hilbert}]{RefinementSpaces}:
$\dim_\R X_n^{0, \bullet} = |V(G_n^\bullet)| \cdot 32$ and
$\dim_\R X_n^{1, \bullet} = |E(G_n^\bullet)| \cdot 7$, both
finite.
\end{proposition}

\begin{proof}
Direct from \cite[Proposition~\texttt{prop:Xn-Hilbert}]{RefinementSpaces}.
\end{proof}

\subsection{Discrete coboundary and adjoint}

\begin{definition}[Coboundary operator]
\label{def:coboundary}
The \emph{discrete coboundary}
$d_n \colon X_n^{0, \bullet} \to X_n^{1, \bullet}$ is defined by
\[
  (d_n f)(e) := \pi_{F^0 \to F^1}\bigl(
    f(\text{head}(e)) - f(\text{tail}(e))\bigr)
\]
for $e \in E^{\rm or}(G_n^\bullet)$ a level-$n$ oriented edge,
where $\text{head}(e), \text{tail}(e) \in V(G_n^\bullet)$ are its
endpoints (with orientation), $f(\text{head}(e)) -
f(\text{tail}(e))$ is computed in the fibre $F^0$ of
\cite[Notation~\texttt{not:fibres}]{RefinementSpaces}, and
$\pi_{F^0 \to F^1} \colon F^0 \to F^1 = \Im(\OO)$ is the
canonical orthogonal projection: $F^0 \to V_8 = \OO$ (component
projection in the orthogonal direct sum
$F^0 = V_{H_4} \oplus V_{D_4} \oplus V_{16} \oplus V_8$),
followed by $\OO \to \Im(\OO)$ (projection onto the imaginary
part).

$d_n$ satisfies the antisymmetry convention on oriented edges:
$(d_n f)(-e) = f(\text{head}(-e)) - f(\text{tail}(-e)) =
f(\text{tail}(e)) - f(\text{head}(e)) = -(d_n f)(e)$, so
$d_n f \in X_n^{1, \bullet}$ (which incorporates the antisymmetry
convention of \cite[Definition~\texttt{def:Xn-1}]{RefinementSpaces}).
\end{definition}

\begin{definition}[Adjoint coboundary]
\label{def:coboundary-adjoint}
The \emph{adjoint coboundary} $d_n^* \colon X_n^{1, \bullet} \to
X_n^{0, \bullet}$ is the $\ell^2$-adjoint of $d_n$ with respect
to the inner products of $X_n^{0, \bullet}$ and $X_n^{1, \bullet}$:
for $A \in X_n^{1, \bullet}$ and $v \in V(G_n^\bullet)$, write
$\iota_{F^1 \to F^0} \colon F^1 = \Im(\OO) \hookrightarrow F^0$
for the canonical inclusion of the imaginary octonion summand
into the $V_8$ summand of $F^0$. Then
\[
  (d_n^* A)(v) = \iota_{F^1 \to F^0}\Bigl(
    \tfrac{1}{2} \sum_{e : v = \text{head}(e)}
      A(e) - \tfrac{1}{2} \sum_{e : v = \text{tail}(e)} A(e)\Bigr),
\]
where each sum runs over oriented edges with $v$ as head or tail
respectively. The factor $\tfrac{1}{2}$ compensates for the
$\tfrac{1}{2}$ in the $X_n^1$ inner product convention
\cite[Definition~\texttt{def:Xn-1}]{RefinementSpaces}; the
$\iota_{F^1 \to F^0}$ inclusion makes $(d_n^* A)(v) \in F^0$ as
required by the codomain $X_n^{0, \bullet}$.
\end{definition}

\begin{lemma}[Coboundary-adjoint is the inner-product adjoint]
\label{lem:dn-adjoint}
$\langle d_n f, A \rangle_{X_n^1} = \langle f, d_n^* A
\rangle_{X_n^0}$ for all $f \in X_n^{0, \bullet}, A \in
X_n^{1, \bullet}$.
\end{lemma}

\begin{proof}
Direct computation. Writing out both sides using
Definitions~\ref{def:coboundary} and
\ref{def:coboundary-adjoint} and the inner product conventions:
the matching of sums over endpoints is standard discrete-Hodge
bookkeeping \cite[\S 1.2]{Chung} adapted to the antisymmetry
convention.
\end{proof}

\subsection{Block operators on $X_n$}

\begin{definition}[Graph Laplacian block $A_n$]
\label{def:An}
The \emph{graph Laplacian block} $A_n \colon X_n \to X_n$ is
\[
  A_n := \begin{pmatrix} d_n^* d_n & 0 \\ 0 & d_n d_n^* \end{pmatrix},
\]
acting on $X_n = X_n^{0, \bullet} \oplus X_n^{1, \bullet}$. The
top-left block $d_n^* d_n$ is the projected graph-Laplacian
block on $X_n^0$ induced by the $\Im(\OO)$ projection in
Definition~\ref{def:coboundary} (positive semi-definite by
construction; it acts as the classical combinatorial graph
Laplacian on the $V_8$-imaginary component of each fibre and
trivially on $V_{H_4}, V_{D_4}, V_{16}$ and the real octonion
direction). The bottom-right block $d_n d_n^*$ is the
edge-Laplacian on $X_n^1$ (also positive semi-definite).
\end{definition}

\begin{definition}[Coboundary block $B_n$]
\label{def:Bn}
The \emph{coboundary block} $B_n \colon X_n \to X_n$ is the
self-adjoint extension of the off-diagonal coboundary-adjoint
pair:
\[
  B_n := \begin{pmatrix} 0 & d_n^* \\ d_n & 0 \end{pmatrix}.
\]
$B_n$ is self-adjoint by Lemma~\ref{lem:dn-adjoint}. Its square is
$B_n^2 = \begin{pmatrix} d_n^* d_n & 0 \\ 0 & d_n d_n^* \end{pmatrix}
= A_n$, so $B_n$ is a ``square root''-type structure for the
Hodge Laplacian.
\end{definition}

\begin{definition}[Mass block $C_n$]
\label{def:Cn}
The \emph{mass block} $C_n \colon X_n \to X_n$ is
\[
  C_n := \begin{pmatrix} \mathbf{P}_{V_{D_4}} & 0 \\ 0 & \id_{X_n^1}
    \end{pmatrix},
\]
where $\mathbf{P}_{V_{D_4}}$ is the orthogonal projection on
$X_n^{0, \bullet}$ onto the $V_{D_4}$-sector component of the
fibre $F^0$ (so $\mathbf{P}_{V_{D_4}}$ acts pointwise on each
vertex value by the linear projection $F^0 = V_{H_4} \oplus
V_{D_4} \oplus V_{16} \oplus V_8 \to V_{D_4}$). $C_n$ is
self-adjoint and positive semi-definite.
\end{definition}

\begin{proposition}[Bounds on the block operators]
\label{prop:block-bounds}
For every $n \geq 0$:
\begin{enumerate}
\item $d_n, d_n^*$ are bounded linear maps with $\|d_n\|_{\op},
  \|d_n^*\|_{\op} \leq 2 \sqrt{\deg(G_n^\bullet)}$, where
  $\deg(G_n^\bullet)$ is the maximum vertex degree of
  $G_n^\bullet$.
\item $A_n$ is bounded self-adjoint with $0 \leq A_n \leq 4
  \deg(G_n^\bullet) \cdot \id$.
\item $B_n$ is bounded self-adjoint with $\|B_n\|_{\op} \leq 2
  \sqrt{\deg(G_n^\bullet)}$.
\item $C_n$ is bounded self-adjoint with $0 \leq C_n \leq \id$.
\end{enumerate}
All four operators are bounded in operator norm uniformly on
$X_n$ (but not uniformly in $n$, since $\deg(G_n^\bullet)$ may
grow with refinement).
\end{proposition}

\begin{proof}
\emph{(1) Coboundary bound.} For $f \in X_n^{0, \bullet}$,
\[
  \|d_n f\|^2_{X_n^1} = \tfrac{1}{2} \sum_e \|f(\text{head}(e)) -
    f(\text{tail}(e))\|^2_{\Im(\OO)}.
\]
By the Cauchy--Schwarz inequality applied to the projection onto
$\Im(\OO)$ and the triangle inequality,
$\|f(\text{head}(e)) - f(\text{tail}(e))\|^2 \leq 2(\|f(\text{head}(e))
\|^2 + \|f(\text{tail}(e))\|^2)$. Summing over $e$:
\[
  \|d_n f\|^2 \leq \sum_e (\|f(\text{head}(e))\|^2 +
    \|f(\text{tail}(e))\|^2) \leq 2 \deg(G_n^\bullet) \cdot
    \|f\|^2,
\]
using that each vertex $v$ appears as head or tail of at most
$\deg(G_n^\bullet)$ edges. Hence $\|d_n\|_{\op} \leq \sqrt{2
\deg(G_n^\bullet)} \leq 2\sqrt{\deg(G_n^\bullet)}$. The same
bound for $d_n^* = (d_n)^*$ follows from $\|T^*\|_{\op} = \|T\|_{\op}$
for bounded operators on Hilbert spaces.

\emph{(2), (3) Laplacian and coboundary block bounds.} $A_n$ is
positive semi-definite as a sum of $T^*T$-type blocks (both
$d_n^* d_n$ and $d_n d_n^*$ are $T^*T$ or $T T^*$). The upper
bound $A_n \leq 4 \deg(G_n^\bullet) \id$ follows from $\|d_n^* d_n
\|_{\op} \leq \|d_n\|_{\op}^2 \leq 4 \deg(G_n^\bullet)$.

For $B_n = \begin{pmatrix} 0 & d_n^* \\ d_n & 0 \end{pmatrix}$:
\[
  \|B_n(f, A)\|^2 = \|d_n^* A\|^2 + \|d_n f\|^2 \leq
    \|d_n\|_{\op}^2 (\|A\|^2 + \|f\|^2) = \|d_n\|_{\op}^2 \|(f, A)\|^2,
\]
so $\|B_n\|_{\op} \leq \|d_n\|_{\op} \leq 2\sqrt{\deg(G_n^\bullet)}$.

\emph{(4) Mass block bound.} $\mathbf{P}_{V_{D_4}}$ is an
orthogonal projection, hence $0 \leq \mathbf{P}_{V_{D_4}} \leq
\id$. The identity on $X_n^1$ is trivially $0 \leq \id \leq \id$.
Combining the blocks: $0 \leq C_n \leq \id$.
\end{proof}

\begin{remark}[Bounds depend on $n$]
\label{rmk:bounds-depend}
The operator bounds in Proposition~\ref{prop:block-bounds} involve
$\deg(G_n^\bullet)$, which grows with $n$ as the refinement
subdivides edges and centroid vertices acquire degree at least
the lengths of their parent face boundaries. This is expected
and does not affect the finite-level theorems of this paper. No
inverse-limit boundedness statement for these operators is
proved here; in particular, P3 constructs only a vertex
inverse-limit Hilbert space
(\cite[Theorem~\texttt{thm:Xinv-Hilbert}]{RefinementSpaces}) and
explicitly does not construct an edge inverse-limit Hilbert
space (\cite[Remark~\texttt{rmk:edge-inv}]{RefinementSpaces}).
\end{remark}

%=====================================================================
\section{Quadratic Closure Functional and Gradient Operator}
\label{sec:quadratic}
%=====================================================================

\subsection{The three quadratic pieces}

\begin{definition}[Curvature, kinetic, mass quadratic forms]
\label{def:RnEnQn}
For $\Phi \in X_n$, define:
\begin{align*}
  R_n(\Phi) &:= \tfrac{1}{2} \langle \Phi, A_n \Phi \rangle_{X_n}
    \quad \text{(curvature / graph-Laplacian form)}, \\
  E_n(\Phi) &:= \tfrac{1}{2} \langle \Phi, B_n \Phi \rangle_{X_n}
    \quad \text{(kinetic / coboundary form)}, \\
  Q_n(\Phi) &:= \tfrac{1}{2} \langle \Phi, C_n \Phi \rangle_{X_n}
    \quad \text{(mass form)}.
\end{align*}
$R_n$ and $Q_n$ are positive semi-definite on $X_n$ (by
positive semi-definiteness of $A_n$ and $C_n$). $E_n$ is
indefinite in general (since $B_n$ is self-adjoint but not
positive semi-definite: it has pairs of eigenvalues
$\pm\sqrt{\lambda}$ for each $\lambda \geq 0$ in
$\spec(d_n^* d_n)$).
\end{definition}

\subsection{The closure functional}

\begin{definition}[Closure functional $F_n$]
\label{def:Fn}
Fix real constants $\alpha, \beta \geq 0$, $\gamma > 0$. The
default working regime, used unless explicitly overridden, is
$\alpha, \beta, \gamma > 0$. The \emph{mass-only specialization}
is the case $\alpha = \beta = 0$, $\gamma > 0$. The
\emph{closure functional at level $n$} is
\[
  F_n(\Phi) := \alpha R_n(\Phi) + \beta E_n(\Phi) - \gamma Q_n(\Phi),
\]
a continuous quadratic form on $X_n$.
\end{definition}

\begin{theorem}[Gradient operator]
\label{thm:gradient-operator}
For every $n \geq 0$:
\begin{enumerate}
\item $F_n$ is a continuous symmetric quadratic form on $X_n$,
  equivalently $F_n(\Phi) = \tfrac{1}{2} \langle \Phi, L_n \Phi
  \rangle_{X_n}$ for a unique bounded self-adjoint operator
  $L_n \in B(X_n)$.
\item The gradient of $F_n$ at $\Phi \in X_n$ is
  $\nabla F_n(\Phi) = L_n \Phi$.
\item $L_n = \alpha A_n + \beta B_n - \gamma C_n$, with
  $\|L_n\|_{\op} \leq \alpha \cdot 4\deg(G_n^\bullet) + \beta \cdot
  2\sqrt{\deg(G_n^\bullet)} + \gamma$.
\end{enumerate}
\end{theorem}

\begin{proof}
\emph{(1) Self-adjoint representation.} By
Proposition~\ref{prop:block-bounds}, $A_n, B_n, C_n$ are each
bounded self-adjoint on $X_n$. Their real linear combination
$L_n := \alpha A_n + \beta B_n - \gamma C_n$ is therefore bounded
self-adjoint. By the finite-dimensional version of the
Riesz representation theorem for quadratic forms
\cite[Prop.~I.3.14]{ReedSimonI}, every continuous symmetric
quadratic form on a finite-dimensional Hilbert space is uniquely
represented as $\tfrac{1}{2} \langle \Phi, L \Phi \rangle$ for a
bounded self-adjoint $L$. The quadratic form $F_n$ is manifestly
continuous (polynomial in matrix entries) and symmetric (each of
$R_n, E_n, Q_n$ is symmetric), so this representation exists. By
uniqueness, $L_n = \alpha A_n + \beta B_n - \gamma C_n$ is
\emph{the} representing operator.

\emph{(2) Gradient.} For any bounded self-adjoint $L$ on a
finite-dimensional Hilbert space, the quadratic form $Q(\Phi) :=
\tfrac{1}{2} \langle \Phi, L \Phi \rangle$ is Fr\'echet
differentiable with derivative $\nabla Q(\Phi) = L \Phi$
\cite[\S 4.6]{HornJohnson}.

\emph{(3) Operator-norm bound.} By the triangle inequality:
\[
  \|L_n\|_{\op} \leq \alpha \|A_n\|_{\op} + \beta \|B_n\|_{\op} +
    \gamma \|C_n\|_{\op} \leq \alpha \cdot 4\deg(G_n^\bullet) +
    \beta \cdot 2\sqrt{\deg(G_n^\bullet)} + \gamma,
\]
using Proposition~\ref{prop:block-bounds}.
\end{proof}

\begin{remark}[Sign of $L_n$]
\label{rmk:Ln-sign}
$L_n = \alpha A_n + \beta B_n - \gamma C_n$ is \emph{not} in
general positive semi-definite: the indefinite term $\beta B_n$
and the subtracted mass term $-\gamma C_n$ can produce negative
eigenvalues. Whether $L_n \geq 0$ depends on the specific
$(\alpha, \beta, \gamma)$ and on spectral properties of the
constituent operators. We will \emph{not} assume $L_n \geq 0$ in
general; the continuous-time flow below is defined and studied
for arbitrary bounded self-adjoint $L_n$, and the
norm-contractivity statements are restricted to coercive
subspaces where positivity holds locally.
\end{remark}

%=====================================================================
\section{Continuous-Time Gradient Flow}
\label{sec:continuous-flow}
%=====================================================================

\begin{theorem}[Gradient flow exists and is unique]
\label{thm:flow-exists}
For every initial datum $\Phi_{n, 0} \in X_n$, the ordinary
differential equation
\begin{equation}
\label{eq:flow}
  \dot \Phi_n(t) = -L_n \Phi_n(t), \qquad \Phi_n(0) = \Phi_{n, 0}
\end{equation}
has a unique global-in-time $C^1$ solution $\Phi_n \colon
[0, \infty) \to X_n$ given by
\[
  \Phi_n(t) = e^{-t L_n} \Phi_{n, 0},
\]
where $e^{-t L_n} := \sum_{k=0}^\infty \frac{(-t)^k L_n^k}{k!}$ is
the finite-dimensional matrix exponential of $-t L_n$. The map
$e^{-t L_n} \colon X_n \to X_n$ is a bounded linear operator with
$\|e^{-t L_n}\|_{\op} \leq e^{t \|L_n\|_{\op}}$ for every $t \geq 0$
(and $\|e^{-t L_n}\|_{\op} \leq e^{|t| \|L_n\|_{\op}}$ for every $t \in
\R$ when the extension to negative times is used;
Cor.~\ref{cor:flow-reversible}), and the family
$(e^{-t L_n})_{t \geq 0}$ is a $C^0$-semigroup on $X_n$:
$e^{0 \cdot L_n} = \id$ and $e^{-(t + s) L_n} = e^{-t L_n}
e^{-s L_n}$ for $t, s \geq 0$.
\end{theorem}

\begin{proof}
$L_n$ is a bounded linear operator on a finite-dimensional
Hilbert space, so the linear ODE \eqref{eq:flow} has a unique
solution for every initial datum, by the standard
Picard--Lindel\"of theorem for linear ODEs with bounded
coefficients \cite[Theorem~1.1]{CoddingtonLevinson}; the solution
is given by the matrix exponential
$\Phi_n(t) = e^{-t L_n} \Phi_{n, 0}$
\cite[\S 6.2]{HornJohnson}. The matrix exponential has operator
norm bounded by $e^{t\|L_n\|_{\op}}$ (direct series estimate); on a
finite-dimensional space the series converges in operator norm.
Semigroup property: by the functional-equation identity
$e^{A} e^{B} = e^{A+B}$ for commuting $A, B$, taking $A = -t L_n$
and $B = -s L_n$ gives $e^{-(t+s) L_n}$.
\end{proof}

\begin{corollary}[The matrix exponential $e^{-tL_n}$ is invertible
on $X_n$ for every $t \in \R$]
\label{cor:flow-reversible}
The map $t \mapsto e^{-t L_n}$ extends to all $t \in \R$ as a
one-parameter group of bounded operators on $X_n$. In particular,
$e^{-t L_n}$ is invertible on $X_n$ for every $t \in \R$.
\end{corollary}

\begin{proof}
$(e^{-t L_n})^{-1} = e^{t L_n}$ by direct computation. The
time-reversed ODE $\dot \Phi_n(s) = L_n \Phi_n(s)$ has the same
well-posedness theorem applied to $-L_n$ in place of $L_n$.
\end{proof}

\begin{remark}[Finite-dimensional semigroup]
\label{rmk:finite-dim-semigroup}
Because $X_n$ is finite-dimensional, the distinction between
``bounded semigroup'' and ``$C_0$-semigroup'' that matters in
infinite-dimensional PDE theory is absent here. The flow
$e^{-t L_n}$ is simply a matrix exponential, and all analytic
properties (continuity in $t$, differentiability, strong/uniform
topology equivalence) follow from finite-dimensional linear
algebra.
\end{remark}

%=====================================================================
\section{Energy Dissipation}
\label{sec:energy}
%=====================================================================

\begin{proposition}[Energy dissipation identity]
\label{prop:energy-dissipation}
For every initial datum $\Phi_{n, 0} \in X_n$ and every $t \geq
0$, with $\Phi_n(t) = e^{-t L_n} \Phi_{n, 0}$:
\[
  \frac{d}{dt} F_n(\Phi_n(t)) = -\|L_n \Phi_n(t)\|^2_{X_n} \leq 0.
\]
Hence $F_n(\Phi_n(t))$ is non-increasing in $t$, and the integrated
dissipation identity holds:
\[
  F_n(\Phi_{n, 0}) - F_n(\Phi_n(t)) = \int_0^t \|L_n \Phi_n(s)
  \|^2_{X_n} \, ds.
\]
\end{proposition}

\begin{proof}
Using the gradient formula $\nabla F_n(\Phi) = L_n \Phi$
(Theorem~\ref{thm:gradient-operator}(2)) and the flow equation
$\dot \Phi_n = -L_n \Phi_n$:
\begin{align*}
  \frac{d}{dt} F_n(\Phi_n(t))
  &= \langle \nabla F_n(\Phi_n(t)), \dot \Phi_n(t) \rangle_{X_n}
  = \langle L_n \Phi_n(t), -L_n \Phi_n(t) \rangle_{X_n} \\
  &= -\|L_n \Phi_n(t)\|^2_{X_n}.
\end{align*}
The sign is $\leq 0$ since $\|\cdot\|^2 \geq 0$. Integrating over
$[0, t]$ gives the integrated identity.
\end{proof}

\begin{remark}[Energy decay vs norm decay: distinct theorems]
\label{rmk:energy-vs-norm}
Proposition~\ref{prop:energy-dissipation} says $F_n$ decreases
along the gradient flow. It does \emph{not} say $\|\Phi_n(t)\|_{X_n}$
decreases. In general these are different statements:
\begin{itemize}
\item Energy monotonicity holds unconditionally (since
  $-\|L_n \Phi_n(t)\|^2 \leq 0$).
\item Norm contractivity $\|\Phi_n(t)\|_{X_n} \leq \|\Phi_n(0)
  \|_{X_n}$ holds iff $e^{-t L_n}$ is a contraction for all $t
  \geq 0$, which requires $L_n \geq 0$
  (Corollary~\ref{cor:norm-decay} below).
\end{itemize}
These are separate hypotheses; conflating them would invalidate
the paper.
\end{remark}

\begin{corollary}[Norm contractivity under positivity]
\label{cor:norm-decay}
If $L_n \geq 0$ (i.e., $L_n$ is positive semi-definite), then
$\|e^{-t L_n}\|_{\op} \leq 1$ for all $t \geq 0$, and hence
$\|\Phi_n(t)\|_{X_n} \leq \|\Phi_{n, 0}\|_{X_n}$ for all $t \geq
0$.
\end{corollary}

\begin{proof}
For self-adjoint $L_n$ with $L_n \geq 0$, the spectral theorem
gives $L_n = \int \lambda \, dE(\lambda)$ with
$\lambda \geq 0$ on the spectral support, so $e^{-t L_n} = \int
e^{-t\lambda} \, dE(\lambda)$ with $|e^{-t\lambda}| \leq 1$ for
$\lambda \geq 0$ and $t \geq 0$. Hence $\|e^{-t L_n}\|_{\op} \leq 1$
\cite[Thm.~VIII.5]{ReedSimonI}.

If $L_n$ has any negative eigenvalue $-\mu < 0$, the corresponding
eigenvector grows as $e^{t\mu}$, violating contractivity.
\end{proof}

%=====================================================================
\section{Discrete-Time Euler Step}
\label{sec:discrete-time}
%=====================================================================

\subsection{The Euler iteration}

\begin{definition}[Euler operator]
\label{def:Euler}
For $n \geq 0$ and step size $\eta > 0$, the \emph{Euler operator}
is
\[
  G_{n, \eta} := I - \eta L_n \in B(X_n).
\]
The \emph{Euler iteration} is the discrete-time dynamics
\[
  \Phi_{n, k+1} := G_{n, \eta} \Phi_{n, k} = (I - \eta L_n)
    \Phi_{n, k}, \qquad k = 0, 1, 2, \ldots,
\]
with initial datum $\Phi_{n, 0} \in X_n$.
\end{definition}

\subsection{Energy monotonicity under small step size}

\begin{proposition}[Euler step is energy-decreasing for small $\eta$]
\label{prop:euler-energy-decrease}
Let $\Phi_{n, k} \in X_n$ with $L_n \Phi_{n, k} \neq 0$, and set
$a := \|L_n \Phi_{n, k}\|^2 > 0$ and $b := \langle L_n \Phi_{n, k},
L_n^2 \Phi_{n, k} \rangle_{X_n} \in \R$. The exact per-step
energy change is
\[
  F_n(\Phi_{n, k+1}) - F_n(\Phi_{n, k}) = -\eta a + \tfrac{1}{2}
    \eta^2 b.
\]
Hence:
\begin{enumerate}
\item If $b \leq 0$, every $\eta > 0$ strictly decreases $F_n$.
\item If $b > 0$, every $\eta$ with $0 < \eta < 2a/b$ strictly
  decreases $F_n$, with the optimum at $\eta^* = a/b$ giving
  $F_n(\Phi_{n, k+1}) - F_n(\Phi_{n, k}) = -a^2/(2b) < 0$.
\end{enumerate}
In particular, there exists $\eta_0(\Phi_{n, k}) > 0$ such that
$F_n(\Phi_{n, k+1}) < F_n(\Phi_{n, k})$ for every $0 < \eta <
\eta_0(\Phi_{n, k})$.
\end{proposition}

\begin{proof}
Direct computation:
\begin{align*}
  F_n(\Phi_{n, k+1})
  &= \tfrac{1}{2} \langle \Phi_{n, k} - \eta L_n \Phi_{n, k},
    L_n (\Phi_{n, k} - \eta L_n \Phi_{n, k}) \rangle \\
  &= \tfrac{1}{2} \langle \Phi_{n, k}, L_n \Phi_{n, k} \rangle
    - \eta \langle L_n \Phi_{n, k}, L_n \Phi_{n, k} \rangle
    + \tfrac{1}{2}\eta^2 \langle L_n \Phi_{n, k}, L_n^2 \Phi_{n, k}
      \rangle \\
  &= F_n(\Phi_{n, k}) - \eta a + \tfrac{1}{2}\eta^2 b.
\end{align*}
The claimed sign analysis follows from completing the square in the
quadratic $-\eta a + \tfrac{1}{2}\eta^2 b$: when $b \leq 0$ the
expression is $\leq -\eta a < 0$ for $\eta > 0$; when $b > 0$ it is
negative iff $0 < \eta < 2a/b$, with minimum at $\eta = a/b$.
\end{proof}

\begin{remark}[$b$ is not always positive]
\label{rmk:b-sign}
Note that $b := \langle L_n \Phi, L_n^2 \Phi \rangle$ is not an
inner-product-type quantity unless $L_n \geq 0$, and can be
negative when $L_n$ has negative eigenvalues on the component of
$L_n \Phi$. When $b \leq 0$, the quadratic correction only helps
the descent; when $b > 0$ the step-size bound $\eta < 2a/b$ is
required. The proposition does \emph{not} assume $L_n \geq 0$.
\end{remark}

\begin{remark}[Spectral version on coercive subspace]
\label{rmk:spectral-version-descent}
If $\Phi_{n, k} \in Y_n$ lies in a coercive invariant subspace with
$m_n \id_{Y_n} \leq L_n|_{Y_n} \leq M_n \id_{Y_n}$, then, setting
$v := L_n \Phi_{n,k} \in Y_n$, we have $a = \|v\|^2$ and
$b = \langle v, L_n v \rangle$, so the spectral bound
$m_n \|v\|^2 \leq \langle v, L_n v \rangle \leq M_n \|v\|^2$ yields
$m_n a \leq b \leq M_n a$. The pointwise step-size bound $0 < \eta
< 2a/b$ therefore holds uniformly on $Y_n$ for any $\eta$ with $0 <
\eta < 2/M_n$. The pointwise optimum $\eta = a/b$ depends on
$\Phi_{n,k}$; for a \emph{uniform} choice that minimises the
worst-case norm-contraction rate one takes $\eta = 2/(m_n + M_n)$
(cf.\ Theorem~\ref{thm:coercive-contraction}), but this is not in
general the pointwise energy-descent optimum.
\end{remark}

\subsection{Norm contractivity on coercive invariant subspaces}

\begin{definition}[Coercive invariant subspace]
\label{def:coercive}
A subspace $Y_n \subseteq X_n$ is \emph{coercive invariant} for
$L_n$ if:
\begin{enumerate}
\item $L_n Y_n \subseteq Y_n$ (invariance), and
\item there exist constants $0 < m_n \leq M_n < \infty$ such that
  $m_n \id_{Y_n} \leq L_n|_{Y_n} \leq M_n \id_{Y_n}$
  (coercivity with spectral gap $m_n > 0$).
\end{enumerate}
\end{definition}

\begin{theorem}[Strict contraction on coercive invariant subspaces]
\label{thm:coercive-contraction}
Suppose $Y_n \subseteq X_n$ is a coercive invariant subspace for
$L_n$ with constants $0 < m_n \leq M_n$. Then for every step size
$\eta$ with $0 < \eta < 2/M_n$, the Euler operator $G_{n, \eta} =
I - \eta L_n$ is a strict contraction on $Y_n$:
\[
  \|G_{n, \eta}|_{Y_n}\|_{\op} \leq \max(|1 - \eta m_n|, |1 - \eta
    M_n|) < 1.
\]
The optimal contraction rate $\|G_{n, \eta}\|_{\op}$ is minimised
at $\eta^* = 2/(m_n + M_n)$, giving
$\|G_{n, \eta^*}|_{Y_n}\|_{\op} = (M_n - m_n)/(M_n + m_n) =: \rho_n
< 1$.
\end{theorem}

\begin{proof}
On $Y_n$, $L_n$ is a self-adjoint operator with spectrum in the
interval $[m_n, M_n]$. The operator $I - \eta L_n$ on $Y_n$ has
spectrum in $[1 - \eta M_n, 1 - \eta m_n]$. For $0 < \eta < 2/M_n$,
both endpoints have $|1 - \eta \lambda| < 1$ iff $0 < \eta < 2/
\lambda$. The operator norm on $Y_n$ equals the maximum of
$|1 - \eta \lambda|$ over $\lambda \in [m_n, M_n]$, which is the
larger of $|1 - \eta m_n|$ and $|1 - \eta M_n|$. The optimum is
achieved when these two are equal in magnitude and opposite in
sign: $1 - \eta m_n = -(1 - \eta M_n)$, giving $\eta = 2/(m_n +
M_n)$ and the stated rate.
\end{proof}

\begin{proposition}[No global contraction without coercivity]
\label{prop:no-global-contraction}
If $L_n$ has \emph{any} negative eigenvalue on $X_n$, then
$G_{n, \eta} = I - \eta L_n$ is \emph{not} a contraction on
$X_n$ for any $\eta > 0$: there exists $\xi \in X_n$ with
$\|G_{n, \eta} \xi\| > \|\xi\|$.
\end{proposition}

\begin{proof}
If $L_n \xi_0 = -\mu \xi_0$ for some unit eigenvector $\xi_0 \in
X_n$ and $\mu > 0$, then $G_{n, \eta} \xi_0 = (1 + \eta\mu)
\xi_0$, which has norm $1 + \eta\mu > 1$.
\end{proof}

\begin{remark}[The error we are not making]
\label{rmk:no-global-contraction-error}
Propositions~\ref{prop:euler-energy-decrease} and
\ref{prop:no-global-contraction} together say: the Euler iteration
can be energy-decreasing (for small $\eta$, on vectors where $L_n$
is active) \emph{and} norm-growing (on vectors in the negative
eigenspace of $L_n$). These are compatible: energy $F_n$ and norm
$\|\Phi\|$ are genuinely different functionals, and $L_n$ can
decrease $F_n$ on one eigenvector while increasing $\|\Phi\|$ on
another. An earlier cascade-correspondence manuscript claimed
$G_{n, \eta}$ is a contraction on $X_n$ for $0 < \eta \leq
\|L_n\|^{-1}$ without any positivity hypothesis; that claim is
false, as shown by Proposition~\ref{prop:no-global-contraction}.
\end{remark}

%=====================================================================
\section{Refinement Compatibility and Symmetry Intertwining}
\label{sec:refinement-compat}
%=====================================================================

\subsection{Refinement compatibility hypothesis}

\begin{definition}[Refinement-compatible operator family]
\label{def:refinement-compat}
A family $(T_n)_{n \geq 0}$ of bounded operators $T_n \in B(X_n)$
is \emph{refinement-compatible} if, for every $n \geq 0$:
\[
  p_{n+1, n} T_{n+1} = T_n p_{n+1, n} \quad \text{on } X_{n+1},
\]
where $p_{n+1, n} \colon X_{n+1} \to X_n$ is the downward bonding
map of \cite[Theorem~\texttt{thm:bonding}]{RefinementSpaces} (extended to the
direct-sum space $X_n = X_n^0 \oplus X_n^1$ in the natural
componentwise way: $p_{n+1, n}(f, A) := (p_{n+1, n}^0 f, p_{n+1, n}
^1 A)$).
\end{definition}

\subsection{Refinement compatibility of the constituent operators}

We verify the coboundary factor relation (Proposition~\ref{prop:coboundary-refinement}, $\tfrac{1}{2}$ factor) and exact mass-block compatibility (Proposition~\ref{prop:adjoint-refinement}); we do \emph{not} prove refinement-compatibility of the full Laplacian block $A_n$ or coboundary block $B_n$ in the $\alpha,\beta>0$ regime, and indeed no such identities are claimed in this paper (see Remark~\ref{rmk:refinement-compat-scope}).

\begin{proposition}[Coboundary refinement-compatibility, with factor $\tfrac{1}{2}$]
\label{prop:coboundary-refinement}
The coboundary family $(d_n)_n$ satisfies
\[
  p_{n+1, n}^1 d_{n+1} = \tfrac{1}{2}\, d_n p_{n+1, n}^0
  \quad \text{on } X_{n+1}^0,
\]
where $p_{n+1, n}^1$ is the parent-fibre averaging bonding map
of \cite[Definition~\texttt{def:p1} (using the edge-parent map of Definition~\texttt{def:edge-parent})]{RefinementSpaces}, $p_{n+1, n}^0$ is the
restriction of \cite[Definition~\texttt{def:p0}]{RefinementSpaces}, and the level-
$(n+1)$ edges subdivide the level-$n$ edges via the canonical
edge-parent map of \cite[Definition~\texttt{def:p1} (using the edge-parent map of Definition~\texttt{def:edge-parent})]{RefinementSpaces}.
\end{proposition}

\begin{proof}
For $f \in X_{n+1}^0$ and $e \in E^{\rm or}(G_n^\bullet)$:
\[
  (d_n (p_{n+1, n}^0 f))(e)
  = (p_{n+1, n}^0 f)(\text{head}(e)) - (p_{n+1, n}^0 f)
    (\text{tail}(e))
  = f(\text{head}(e)) - f(\text{tail}(e))
\]
(using $V(G_n^\bullet) \subset V(G_{n+1}^\bullet)$ and
$p^0$ = restriction from \cite[Definition~\texttt{def:p0}]{RefinementSpaces}).

On the other side, the level-$n$ edge $e$ subdivides at level
$n+1$ into $e_1, e_2$ with $\text{head}(e_1) = \text{tail}(e_2)
= m$ (the midpoint of $e$), so
\begin{align*}
  (d_{n+1} f)(e_1) &= f(m) - f(\text{tail}(e)), \\
  (d_{n+1} f)(e_2) &= f(\text{head}(e)) - f(m).
\end{align*}
The parent-fibre average of $d_{n+1} f$ from
\cite[Definition~\texttt{def:p1} (using the edge-parent map of Definition~\texttt{def:edge-parent})]{RefinementSpaces} is
\[
  (p_{n+1, n}^1 d_{n+1} f)(e)
  = \tfrac{1}{2}\bigl((d_{n+1} f)(e_1) + (d_{n+1} f)(e_2)\bigr)
  = \tfrac{1}{2}\bigl(f(\text{head}(e)) - f(\text{tail}(e))\bigr)
  = \tfrac{1}{2} (d_n p_{n+1, n}^0 f)(e),
\]
which proves the claimed factor-$\tfrac{1}{2}$ identity.
\end{proof}

\begin{remark}[The factor-$\tfrac{1}{2}$ is structural]
\label{rmk:factor-structural}
The factor of $\tfrac{1}{2}$ in
Proposition~\ref{prop:coboundary-refinement} is intrinsic to the
combination of an \emph{undivided-difference coboundary} $d_n$ at
every refinement level with the \emph{averaging bonding map}
$p_{n+1, n}^1$ of \cite[Definition~\texttt{def:p1} (using the edge-parent map of Definition~\texttt{def:edge-parent})]{RefinementSpaces}. Changing the
coboundary normalization (e.g.\ scaling $d_n$ by $2^n$ or by an
edge-length factor) can convert the factor into different constants,
but cannot eliminate it without sacrificing some other structural
property (e.g.\ $\|p_{n+1, n}^1\|_{\op} \leq 1$). We keep the factor
of $\tfrac{1}{2}$ explicit and propagate it through the downstream
operator identities.
\end{remark}

\begin{proposition}[Mass-block refinement compatibility]
\label{prop:adjoint-refinement}
Under the conventions of \S\ref{sec:operators} and
\cite{RefinementSpaces}, the mass block satisfies
\[
  p_{n+1, n} C_{n+1} = C_n p_{n+1, n}
\]
exactly (no scaling factor). The analogous exact identity for the
coboundary, adjoint coboundary, or graph-Laplacian blocks is not
established in this paper; see
Remark~\ref{rmk:refinement-compat-scope} below for the scope of
what is and is not proved.
\end{proposition}

\begin{proof}
$C_n = \begin{pmatrix} \mathbf{P}_{V_{D_4}} & 0 \\ 0 & \id
\end{pmatrix}$ acts pointwise on each vertex / edge value by a
fixed linear operation on fibres: on $X_n^0$ by the projection
$\mathbf{P}_{V_{D_4}} \colon F^0 \to V_{D_4}$, and on $X_n^1$ by
the identity. Both of these commute with fibrewise restriction
(vertex case, $p_{n+1, n}^0$) and with fibrewise parent-fibre
averaging (edge case, $p_{n+1, n}^1$). Hence $p_{n+1, n} C_{n+1} =
C_n p_{n+1, n}$ exactly.
\end{proof}

\begin{remark}[Scope of refinement compatibility]
\label{rmk:refinement-compat-scope}
What is actually proved here:
\begin{itemize}
\item The mass block $C_n$ is refinement-compatible exactly:
  $p C_{n+1} = C_n p$ (Proposition~\ref{prop:adjoint-refinement}).
\item The coboundary obeys the factor-$\tfrac{1}{2}$ relation
  $p^1 d_{n+1} = \tfrac{1}{2} d_n p^0$
  (Proposition~\ref{prop:coboundary-refinement}).
\end{itemize}
What is \emph{not} proved here: an analogous factor-$\tfrac{1}{2}$
relation for $d_n^*$; a block identity of the form $p B_{n+1} = c
B_n p$ for any explicit constant $c$; an induced factor on the
graph-Laplacian block $A_n$. Each of these would require a separate
calculation reconciling the distinct $p^0$ and $p^1$ bonding-map
factors across the off-diagonal blocks, and none is carried out in
this paper.

Consequently, we restrict Theorem~\ref{thm:refinement-compat}
(and Theorem~\ref{thm:flow-intertwining} below) to the mass-only
regime $\alpha = \beta = 0$, where refinement compatibility is
positively established. For arbitrary $(\alpha, \beta, \gamma)$ we
make \emph{no} claim about refinement-intertwining of $L_n$: neither
that it holds (no such calculation is offered), nor that it fails
(no negative example is offered). Downstream papers requiring
refinement-intertwining outside the mass-only regime must supply
their own analysis of the coboundary--adjoint block identities.
\end{remark}

\subsection{The refinement-intertwining theorem: mass-only case}

\begin{theorem}[$L_n$ is refinement-compatible in the mass-only case]
\label{thm:refinement-compat}
Under the restriction $\alpha = \beta = 0$ (i.e., $L_n = -\gamma
C_n$), the gradient operator family $(L_n)_{n \geq 0}$ is
refinement-compatible:
\[
  p_{n+1, n} L_{n+1} = L_n p_{n+1, n}
  \quad \text{for every } n \geq 0.
\]
For arbitrary $(\alpha, \beta, \gamma)$ we make no claim about the
corresponding identity; see
Remark~\ref{rmk:refinement-compat-scope}.
\end{theorem}

\begin{proof}
Under $\alpha = \beta = 0$, $L_n = -\gamma C_n$. By
Proposition~\ref{prop:adjoint-refinement}, $p_{n+1, n} C_{n+1} =
C_n p_{n+1, n}$, hence (multiplying by $-\gamma$) $p_{n+1, n}
L_{n+1} = L_n p_{n+1, n}$.
\end{proof}

\begin{theorem}[Flow intertwining: mass-only case]
\label{thm:flow-intertwining}
Under the restriction $\alpha = \beta = 0$, the continuous-time
flow of $L_n$ intertwines with the bonding maps:
\[
  p_{n+1, n} e^{-t L_{n+1}} = e^{-t L_n} p_{n+1, n}
  \quad \text{for every } t \geq 0.
\]
For arbitrary $(\alpha, \beta, \gamma)$ we make no claim about the
corresponding identity; see
Remark~\ref{rmk:refinement-compat-scope}.
\end{theorem}

\begin{proof}
Under $\alpha = \beta = 0$, Theorem~\ref{thm:refinement-compat}
gives $p_{n+1, n} L_{n+1} = L_n p_{n+1, n}$ exactly. By functional
calculus for bounded operators on finite-dimensional Hilbert
spaces \cite[Thm.~VIII.5]{ReedSimonI}, applied to the power series
$e^{-tx} = \sum_{k \geq 0} (-tx)^k / k!$ term by term (each power
$L_n^k$ commutes with $p_{n+1, n}$ by iteration of the linear
intertwining), we obtain
$p_{n+1, n} e^{-t L_{n+1}} = e^{-t L_n} p_{n+1, n}$.
\end{proof}

\begin{remark}[Pedagogical note on §7 scope]
\label{rmk:section7-scope}
Section~\ref{sec:refinement-compat} \emph{does not} establish
refinement-intertwining of the full cascade gradient flow with
arbitrary $\alpha, \beta, \gamma$. It establishes:
\begin{itemize}
\item an explicit factor-$\tfrac{1}{2}$ relation for the
  coboundary (Proposition~\ref{prop:coboundary-refinement});
\item exact compatibility for the mass block
  (Proposition~\ref{prop:adjoint-refinement});
\item exact compatibility for $L_n$ and its flow in the mass-only
  degenerate case
  (Theorems~\ref{thm:refinement-compat}, \ref{thm:flow-intertwining}).
\end{itemize}
Any downstream paper needing a refinement-intertwining theorem
for non-degenerate $(\alpha, \beta, \gamma)$ must supply its own
analysis. This paper makes no claim about whether such intertwining
holds in that regime; it is simply not established here.
\end{remark}

\subsection{Symmetry intertwining}

\begin{corollary}[Symmetries commuting with $L_n$ commute with the flow]
\label{cor:symmetry-flow}
If $g \in W^\bullet$ acts on $X_n$ by $\rho_n(g) \in B(X_n)$ with
$\rho_n(g) L_n = L_n \rho_n(g)$ (i.e., $g$-equivariant gradient
operator), then $\rho_n(g) e^{-t L_n} = e^{-t L_n} \rho_n(g)$ for
every $t \geq 0$.
\end{corollary}

\begin{proof}
Same functional-calculus argument as
Theorem~\ref{thm:flow-intertwining}.
\end{proof}

%=====================================================================
\section{Sectorwise Intertwiners: Explicit Bounds}
\label{sec:intertwiners}
%=====================================================================

\subsection{Admissible finite-level intertwiners}

\begin{definition}[Generator list]
\label{def:generators}
The \emph{finite-level generator list} $\mathcal{G}_n$ at level
$n$ consists of the following bounded real-Hilbert-space operators:
\begin{itemize}
\item downward bonding maps: $p_{n+1, n}^0, p_{n+1, n}^1$ and
  their composition $p_{n+1, n}$ on $X_{n+1} \to X_n$;
\item upward cylindrical embeddings: $j_{n, n+1}^0, j_{n, n+1}^1$
  and $j_{n, n+1}$ on $X_n \to X_{n+1}$;
\item sector projections: $\mathbf{P}_{V_{H_4}},
  \mathbf{P}_{V_{D_4}}, \mathbf{P}_{V_{16}}, \mathbf{P}_{V_8}$ on
  $X_n^0$;
\item Coxeter actions: $\rho_n(g)$ for $g \in W^\bullet$;
\item discrete differential operators: $d_n, d_n^*, A_n, B_n, C_n,
  L_n$ from \S\ref{sec:operators};
\item flow-operator family: $\{e^{-t L_n} : t \geq 0\}$;
\item the $G_2$-stabilizer fibre map $\Im(\OO) \to \Im(\OO)$
  fixing a chosen unit $i_0$ (an orthogonal map of norm $1$,
  imported from \cite{AlgebraicSubstrate}).
\end{itemize}
The rationality-side $\sigma$-involution and the icosian star-map
are \emph{not} elements of $\mathcal{G}_n$: by P3
(\cite[Definition~\texttt{def:sigma}, Theorem~\texttt{thm:commute}]
{RefinementSpaces}) these are defined only on the mixed-form
$\Q(\varphi)/\Q$ vertex space and do not extend to bounded
operators on the real Hilbert space $X_n^0$. They appear in
this paper only at the rational-form level, never as elements
of $\mathcal{G}_n$.
\end{definition}

\begin{definition}[Admissible intertwiner, syntactic]
\label{def:admissible}
An \emph{admissible finite-level intertwiner} is a formal
expression obtained from the generator list $\mathcal{G}_n$ (or
$\mathcal{G}_n \cup \mathcal{G}_m$ for inter-level intertwiners)
by the syntactic rules:
\begin{itemize}
\item Every generator $T \in \mathcal{G}_n$ is an admissible
  intertwiner (of level $n$, or between appropriate levels).
\item If $T_1, T_2$ are admissible intertwiners with compatible
  domains and codomains (matching refinement levels and state-space
  types), then $c_1 T_1 + c_2 T_2$ is an admissible intertwiner
  for any $c_1, c_2 \in \R$.
\item If $T_1 \colon X_m \to X_\ell$ and $T_2 \colon X_n \to X_m$
  are admissible, then the composition $T_1 \circ T_2 \colon X_n
  \to X_\ell$ is admissible.
\end{itemize}
The definition is \emph{syntactic}: it does not assume that the
resulting operator is bounded \emph{a priori}. Boundedness is
proved in Theorem~\ref{thm:intertwiner-bounds} below from the
printed generator bounds.
\end{definition}

\subsection{Explicit bounds for the named generators}

\begin{proposition}[Generator bounds table]
\label{prop:generator-bounds}
The named generators from Definition~\ref{def:generators} satisfy
the following explicit operator-norm bounds on their respective
finite-level Hilbert spaces:
\begin{center}
\begin{tabular}{lll}
\hline
Generator & Norm bound & Source \\
\hline
$p_{n+1, n}^0$ (vertex bonding) & $\|p^0\|_{\op} \leq 1$ &
  \cite[Theorem~\texttt{thm:bonding}]{RefinementSpaces} \\
$p_{n+1, n}^1$ (edge bonding) & $\|p^1\|_{\op} \leq 1$ &
  \cite[Theorem~\texttt{thm:bonding}]{RefinementSpaces} \\
$p_{n+1, n}$ (direct-sum bonding $X_{n+1} \to X_n$) & $\|p\|_{\op}
  \leq 1$ & direct-sum norm, from $\|p^0\|, \|p^1\| \leq 1$ \\
$j_{n, n+1}^0$ (vertex embedding) & $\|j^0\|_{\op} = 1$ &
  \cite[Proposition~\texttt{prop:adjoints}]{RefinementSpaces} (adjoint of restriction,
  isometric on image) \\
$j_{n, n+1}^1$ (edge embedding) & $\|j^1\|_{\op} = 2^{-1/2}$ &
  computation from
  \cite[Proposition~\texttt{prop:adjoints}]{RefinementSpaces}: the
  half-section identity $p^1 j^1 = \tfrac{1}{2} \id$ together with
  unitarity of the inner-product adjoint gives
  $\|j^1 B\|^2 = \langle j^1 B, j^1 B \rangle =
  \langle B, p^1 j^1 B \rangle = \tfrac{1}{2}\|B\|^2$ \\
$j_{n, n+1}$ (direct-sum embedding $X_n \to X_{n+1}$) & $\|j\|_{\op}
  = 1$ & direct-sum norm: $\max(\|j^0\|, \|j^1\|) = 1$ \\
$\mathbf{P}_{V_\bullet}$ (sector projections) & $\|\mathbf{P}\|_{\op}
  = 1$ & orthogonal projections \\
$\rho_n(g)$ (Coxeter action) & $\|\rho_n(g)\|_{\op} = 1$ & unitary,
  \cite[Lemma~\texttt{lem:coxeter-unitary}]{RefinementSpaces} \\
$d_n$ (coboundary) & $\|d_n\|_{\op} \leq 2\sqrt{\deg(G_n^\bullet)}$
  & Prop.~\ref{prop:block-bounds}(1) \\
$d_n^*$ (adjoint coboundary) & $\|d_n^*\|_{\op} \leq 2\sqrt{\deg(G_n
  ^\bullet)}$ & adjoint of $d_n$ \\
$A_n$ (graph-Laplacian block) & $\|A_n\|_{\op} \leq 4\deg(G_n
  ^\bullet)$ & Prop.~\ref{prop:block-bounds}(2) \\
$B_n$ (coboundary block) & $\|B_n\|_{\op} \leq 2\sqrt{\deg(G_n
  ^\bullet)}$ & Prop.~\ref{prop:block-bounds}(3) \\
$C_n$ (mass block) & $0 \leq C_n \leq \id$, so $\|C_n\|_{\op} \leq 1$
  & Prop.~\ref{prop:block-bounds}(4) \\
$L_n$ (gradient operator) & $\|L_n\|_{\op} \leq
  4\alpha\,\deg(G_n^\bullet) + 2\beta\sqrt{\deg(G_n^\bullet)}
  + \gamma$ & Thm.~\ref{thm:gradient-operator}(3) \\
$e^{-t L_n}$ (flow, $t \geq 0$) & $\|e^{-t L_n}\|_{\op} \leq
  e^{t \|L_n\|_{\op}}$ & Thm.~\ref{thm:flow-exists} \\
\hline
\end{tabular}
\end{center}
The $G_2$-stabilizer fibre map (from
\cite[Definitions~\texttt{def:G2}, \texttt{def:stab-u} and
Fact~\texttt{fact:SU3-stab}]{AlgebraicSubstrate}) acts on
$\Im(\OO) = F^1$ as a compact-group element of orthogonal norm $1$;
its $\ell^2$-extension to $X_n^1$ therefore has operator norm
exactly $1$.
\end{proposition}

\begin{proof}
Each row of the table is verified by citation or by the referenced
proposition. $\|\rho_n(g)\|_{\op} = 1$ because $\rho_n(g)$ is
unitary
(\cite[Lemma~\texttt{lem:coxeter-unitary}]{RefinementSpaces}: it
permutes an $\ell^2$-orthonormal basis up to fibre-level
orthogonal transformations). The $G_2$-stabilizer claim uses that
the relevant element is in the compact orthogonal group of
$\Im(\OO)$.
\end{proof}

\begin{theorem}[Admissible intertwiners are bounded, continuous, Borel]
\label{thm:intertwiner-bounds}
Every admissible finite-level intertwiner $T$ (in the syntactic
sense of Definition~\ref{def:admissible}) is a bounded linear
operator on the finite-dimensional Hilbert space in its codomain.
Specifically:
\begin{enumerate}
\item (Boundedness with printed bound.) If $T$ is expressed by
  the syntax of Definition~\ref{def:admissible} as a finite
  sum-of-products of generators
  $T = \sum_i c_i \, T_{i, k_i^1} \circ \cdots \circ T_{i, k_i^{m_i}}$
  with $c_i \in \R$ and each $T_{i, k_i^j}$ in the generator
  list of Proposition~\ref{prop:generator-bounds}, then
  $\|T\|_{\op} \leq \sum_i |c_i| \prod_j \|T_{i, k_i^j}\|_{\op}$,
  using the triangle inequality on sums
  ($\|T_1 + T_2\|_{\op} \leq \|T_1\|_{\op} + \|T_2\|_{\op}$),
  the absolute scalar identity $\|c T\|_{\op} = |c| \|T\|_{\op}$,
  and sub-multiplicativity on compositions
  ($\|T_1 T_2\|_{\op} \leq \|T_1\|_{\op} \|T_2\|_{\op}$).
\item (Continuity.) $T$ is continuous as a linear map between
  finite-dimensional Hilbert spaces (automatic for bounded linear
  maps).
\item (Borel measurability.) $T$ is Borel-measurable (continuity
  plus finite-dimensionality).
\item (Symmetry intertwining, conditional.) Suppose
  $T = \sum_i c_i \, T_{i, k_i^1} \circ \cdots \circ
  T_{i, k_i^{m_i}}$ as in (1), and a fixed symmetry
  $g \in W^\bullet$ satisfies a prescribed intertwining relation
  $\rho(g) T_{i, k_i^j} = T_{i, k_i^j}' \rho(g)$ for every
  generator factor (with the same source-side $\rho(g)$ in every
  case). Then $\rho(g) T = T' \rho(g)$ where
  $T' := \sum_i c_i \, T_{i, k_i^1}' \circ \cdots \circ
  T_{i, k_i^{m_i}}'$, by linearity of $\rho(g)$ on each summand.
\end{enumerate}
\end{theorem}

\begin{proof}
\emph{(1) Boundedness with printed bound.} By structural induction
on the formal expression of $T$ per
Definition~\ref{def:admissible}, using the sub-multiplicativity
of $\|\cdot\|_{\op}$ on compositions and the triangle inequality on
sums. Base case: each generator has a norm bound from
Proposition~\ref{prop:generator-bounds}. Inductive cases (sum,
composition): standard operator-theoretic inequalities
\cite[Ch.~I.1]{ReedSimonI}.

\emph{(2) Continuity.} Bounded linear maps on normed spaces are
continuous \cite[Thm.~I.1.3]{ReedSimonI}.

\emph{(3) Borel measurability.} Continuous maps between
topological spaces are Borel measurable
\cite[Cor.~1.2.12]{Bogachev}. In finite dimension, all reasonable
measurable structures coincide.

\emph{(4) Symmetry intertwining.} For each summand
$c_i \, T_{i, k_i^1} \circ \cdots \circ T_{i, k_i^{m_i}}$,
iterating the generator intertwinings along the composition
gives
$\rho(g) (c_i \, T_{i, k_i^1} \circ \cdots \circ T_{i, k_i^{m_i}})
= c_i (T_{i, k_i^1}' \circ \cdots \circ T_{i, k_i^{m_i}}')
\rho(g)$.
Summing over $i$ and using $\R$-linearity of $\rho(g)$ yields
$\rho(g) T = T' \rho(g)$.
\end{proof}

\begin{remark}[What this replaces]
\label{rmk:intertwiner-replacement}
Theorem~\ref{thm:intertwiner-bounds} is the structural replacement
for the false global theorem ``every correspondence arising from
closure dynamics is Borel and equivariant'' that appeared in an
earlier foundations draft. Instead of one grand theorem, we have:
\begin{itemize}
\item a printed generator list
  (Definition~\ref{def:generators});
\item a printed bounds table
  (Proposition~\ref{prop:generator-bounds});
\item a compositionality theorem
  (Theorem~\ref{thm:intertwiner-bounds}).
\end{itemize}
Downstream papers naming an intertwiner as a specific composition
of generators read off the operator bound, continuity, and
intertwining properties from this table and theorem. No abstract
``Borel/equivariance'' theorem is invoked; every claim is
composable from printed ingredients.
\end{remark}

%=====================================================================
\section{Citation Contract and Downstream Handoff}
\label{sec:handoff}
%=====================================================================

\subsection{What downstream papers may cite}

\begin{enumerate}
\item Definition~\ref{def:Fn} and
  Theorem~\ref{thm:gradient-operator}: the closure functional
  $F_n$ and its gradient operator $L_n = \alpha A_n + \beta B_n -
  \gamma C_n$, a bounded self-adjoint operator on $X_n$.
\item Theorem~\ref{thm:flow-exists}: global existence and
  uniqueness of the continuous-time gradient flow
  $\Phi_n(t) = e^{-t L_n} \Phi_n(0)$.
\item Proposition~\ref{prop:energy-dissipation}: energy
  dissipation along the flow.
\item Propositions~\ref{prop:euler-energy-decrease},
  \ref{prop:no-global-contraction} and
  Theorem~\ref{thm:coercive-contraction}: discrete-time Euler
  analysis, with energy monotonicity and norm contractivity
  cleanly separated.
\item Proposition~\ref{prop:coboundary-refinement}: the explicit
  factor-$\tfrac{1}{2}$ coboundary relation $p^1 d_{n+1} =
  \tfrac{1}{2} d_n p^0$. Proposition~\ref{prop:adjoint-refinement}:
  exact mass-block compatibility $p C_{n+1} = C_n p$.
  Theorems~\ref{thm:refinement-compat} and
  \ref{thm:flow-intertwining}: refinement compatibility of $L_n$
  and of the continuous-time flow \emph{in the mass-only case}
  $\alpha = \beta = 0$. For arbitrary $(\alpha, \beta, \gamma)$
  this paper makes no claim about refinement-intertwining of $L_n$;
  see Remark~\ref{rmk:refinement-compat-scope}.
\item Definitions~\ref{def:generators}, \ref{def:admissible} and
  Theorem~\ref{thm:intertwiner-bounds}: explicit generator-family list
  and bounded-intertwiner theorem replacing the false global
  Borel/equivariance theorem.
\end{enumerate}

\subsection{Recommended citation phrasing}

``By P4 \cite{ClosureDynamics}, each finite refinement level
$n$ carries an explicit quadratic closure functional
$F_n = \tfrac{1}{2}\langle \cdot, L_n \cdot \rangle$ with bounded
self-adjoint gradient $L_n$ (Theorem~\ref{thm:gradient-operator}),
continuous-time semigroup flow $e^{-t L_n}$
(Theorem~\ref{thm:flow-exists}), energy monotonicity
(Proposition~\ref{prop:energy-dissipation}), strict contraction
only on coercive invariant subspaces
(Theorem~\ref{thm:coercive-contraction}), an explicit
factor-$\tfrac{1}{2}$ coboundary refinement relation
(Proposition~\ref{prop:coboundary-refinement}) and exact
mass-block refinement compatibility
(Proposition~\ref{prop:adjoint-refinement}), with full
refinement-intertwining $p_{n+1, n} e^{-t L_{n+1}} = e^{-t L_n}
p_{n+1, n}$ holding in the mass-only case $\alpha = \beta = 0$
(Theorem~\ref{thm:flow-intertwining}); and explicit bounded
intertwiner bounds from the explicit generator-family list
(Theorem~\ref{thm:intertwiner-bounds}).''

\subsection{What this paper does not provide}

\begin{enumerate}
\item No continuum-limit statement about $(L_n)_n$ as $n \to
  \infty$ or about $(e^{-t L_n})_n$.
\item No target-side metric / hydrodynamic / gauge / spectral /
  Hecke / cohomology / CTM/TM correspondence theorem.
\item No PDE identification of $\dot\Phi_n = -L_n \Phi_n$ with
  any classical evolution equation.
\item No universal Borel/equivariance principle beyond the
  explicit generator-family list.
\end{enumerate}

\appendix

%=====================================================================
\section{Notation Audit}
\label{app:notation}
%=====================================================================

\begin{center}
\begin{tabular}{lll}
\hline
Symbol & Meaning & Defined \\
\hline
$X_n$ & combined state space $X_n^0 \oplus X_n^1$ & Def.~\ref{def:Xn} \\
$D_n$ & domain of the operators ($= X_n$) & Def.~\ref{def:Xn} \\
$d_n, d_n^*$ & discrete coboundary and adjoint & Defs.~\ref{def:coboundary}, \ref{def:coboundary-adjoint} \\
$A_n$ & graph-Laplacian block operator & Def.~\ref{def:An} \\
$B_n$ & coboundary block operator & Def.~\ref{def:Bn} \\
$C_n$ & mass block operator & Def.~\ref{def:Cn} \\
$R_n, E_n, Q_n$ & curvature, kinetic, mass quadratic forms & Def.~\ref{def:RnEnQn} \\
$F_n$ & closure functional $\alpha R_n + \beta E_n - \gamma Q_n$ & Def.~\ref{def:Fn} \\
$L_n$ & gradient operator $\alpha A_n + \beta B_n - \gamma C_n$ & Thm.~\ref{thm:gradient-operator}(3) \\
$e^{-t L_n}$ & continuous-time flow & Thm.~\ref{thm:flow-exists} \\
$G_{n, \eta}$ & Euler operator $I - \eta L_n$ & Def.~\ref{def:Euler} \\
$Y_n$ & coercive invariant subspace & Def.~\ref{def:coercive} \\
$m_n, M_n$ & coercive spectral bounds & Def.~\ref{def:coercive} \\
$\mathcal{G}_n$ & finite-level generator list & Def.~\ref{def:generators} \\
\hline
\end{tabular}
\end{center}

\begin{thebibliography}{99}

\bibitem{AlgebraicSubstrate}
\emph{Cascade--Classical Correspondence II: The Cascade Algebraic
Substrate},
\texttt{papers/cascade-algebraic-substrate/cascade-algebraic-substrate.tex}.
(Paper P2.) Pinned results used here, cited by label:
\texttt{def:icosian-star} (icosian star-map $\ast$ on the icosian
ring $I \subset \HH$, identified at the rational level with the
coefficientwise $\sigma$ of P1; this map is used in this paper
only at the rational-form level, never as a real-Hilbert-space
operator on $X_n^0$);
\texttt{def:G2}, \texttt{def:stab-u}, and
Fact~\texttt{fact:SU3-stab} ($G_2$ stabilizer of a unit imaginary
octonion $\cong \mathrm{SU}(3)$);
\texttt{def:Cl-basis}, \texttt{def:octonions}, and
\texttt{def:Q-forms} (sector-fibre $\Q$-forms used in P3's
mixed-form construction).

\bibitem{RefinementSpaces}
\emph{Cascade--Classical Correspondence III: Refinement Spaces
and Symmetry Actions},
\texttt{papers/cascade-refinement-spaces/cascade-refinement-spaces.tex}.
(Paper P3.) Pinned results used here, cited by label:
Notation~\texttt{not:fibres} (sector fibres $F^0, F^1$);
Definitions~\texttt{def:Xn-0} and~\texttt{def:Xn-1} (finite-level
$\ell^2$ spaces $X_n^0, X_n^1$);
Proposition~\texttt{prop:Xn-Hilbert} (finite-level Hilbert
structure);
Definition~\texttt{def:edge-parent} (canonical edge-parent map
on subdividing edges);
Definition~\texttt{def:p0} (vertex bonding $p^0$ as restriction);
Definition~\texttt{def:p1} (edge bonding $p^1$ as parent-fibre
averaging);
Theorem~\texttt{thm:bonding} (operator-norm contraction bounds
for $p^0, p^1$);
Proposition~\texttt{prop:adjoints} (adjoints
$j_{n, n+1} = (p_{n+1, n})^*$, giving
$\|j^0\|_{\op} = 1$ and $\|j^1 B\|^2 = \tfrac{1}{2}\|B\|^2$);
Definition~\texttt{def:coxeter} (branchwise Coxeter action
$\rho_n(g)$);
Lemma~\texttt{lem:coxeter-unitary} (unitarity of $\rho_n(g)$ on
each $X_n^{0, \bullet}, X_n^{1, \bullet}$);
Theorem~\texttt{thm:Xinv-Hilbert} (vertex inverse-limit Hilbert
space; no edge inverse-limit Hilbert space is constructed in
P3, see Remark~\texttt{rmk:edge-inv} of P3).

\bibitem{ClosureDynamics}
(This paper.) \emph{Cascade--Classical Correspondence IV: Discrete
Closure Dynamics}.

\bibitem{ReedSimonI}
M.~Reed and B.~Simon, \emph{Methods of Modern Mathematical
Physics, Vol.~I: Functional Analysis}, rev.\ ed., Academic Press,
1980. Quadratic forms and self-adjoint operators: Ch.~I.3;
bounded operators: Ch.~I.1; spectral theorem and functional
calculus: Ch.~VIII.5.

\bibitem{HornJohnson}
R.~A.~Horn and C.~R.~Johnson, \emph{Matrix Analysis}, 2nd ed.,
Cambridge University Press, 2013. Matrix exponentials and linear
ODEs: Ch.~6.2; spectral decomposition in finite dimension:
Ch.~2.4; Fr\'echet differentiability of quadratic forms on
finite-dimensional Hilbert spaces: \S 4.6 (treated in the chapter
on matrix-valued functions).

\bibitem{Chung}
F.~R.~K.~Chung, \emph{Spectral Graph Theory}, CBMS Regional
Conference Series in Mathematics 92, AMS, 1997. Graph Laplacians,
coboundaries, and finite-graph Hodge theory: Ch.~1--2.

\bibitem{CoddingtonLevinson}
E.~A.~Coddington and N.~Levinson, \emph{Theory of Ordinary
Differential Equations}, McGraw-Hill, 1955. Linear ODEs with
bounded coefficients and global existence: Ch.~1.

\bibitem{Hartman}
P.~Hartman, \emph{Ordinary Differential Equations}, 2nd ed.,
Birkh\"auser, 1982. Linear ODEs: Ch.~IV.

\bibitem{Bogachev}
V.~I.~Bogachev, \emph{Measure Theory}, Vol.~I, Springer, 2007.
Borel measurability of continuous maps: Cor.~1.2.12.

\end{thebibliography}

\end{document}
